\documentclass{atscv}
\usepackage{atscv-cv-support-burgundy-consulting-de}
\begin{document}
\cvdebugoptions
\cvname{Jonas Richter}
\cvtitle{Support- und Service-Operations Lead}
\cvemail{jonas.richter@example.com}
\cvlocation{Hamburg | Remote-first}
\makecvheader

\cvsection{\headingSummary}
\cvprofile{Service-orientierter Operations-Spezialist mit Erfahrung im Aufbau skalierbarer Support-Prozesse für B2B-Softwareprodukte. Schwerpunkt auf Ticket-Qualität, Wissensmanagement und enger Abstimmung zwischen Support, Engineering und Customer Success.}

\cvsection{\headingExperience}
\cventry{Head of Support Operations}{2023--Heute}{Clearline Platforms}{\cvachievement{Führte ein neues Routing- und Priorisierungsmodell ein, reduzierte die Erstreaktionszeit um 46\% und erhöhte die Lösungsquote im First Level signifikant.}}
\cventry{Senior Technical Support Manager}{2020--2023}{Signalwerk Solutions}{\cvachievement{Etablierte KPIs für SLAs und Qualitätsreviews, inklusive strukturiertem Post-Incident-Lernen mit monatlichen Verbesserungsmaßnahmen über mehrere Teams hinweg.}}
\cventry{Technical Support Engineer}{2017--2020}{Mosaic Data Systems}{\cvachievement{Bearbeitete komplexe Eskalationen in API- und Integrationsszenarien und baute eine interne Wissensdatenbank auf, die Einarbeitungszeiten deutlich verkürzte.}}

\cvsection{\headingSkills}
\cvskillgroup{Support-Prozesse}{ITIL-nahe Abläufe, Eskalationsmanagement, SLA-Steuerung, Quality Audits}
\cvskillgroup{Plattformen \& Tools}{Zendesk, Jira, Confluence, SQL für Ticketanalyse, API-Debugging}
\cvskillgroup{Führung \& Zusammenarbeit}{Coaching, Schichtplanung, bereichsübergreifende Kommunikation, Stakeholder-Reporting}
\end{document}
